% Options for packages loaded elsewhere
\PassOptionsToPackage{unicode}{hyperref}
\PassOptionsToPackage{hyphens}{url}
\PassOptionsToPackage{dvipsnames,svgnames,x11names}{xcolor}
\PassOptionsToPackage{space}{xeCJK}
\documentclass[
  a4paper,
  10pt]{article}
\usepackage{xcolor}
\usepackage{amsmath,amssymb}
\setcounter{secnumdepth}{5}
\usepackage{iftex}
\ifPDFTeX
  \usepackage[T1]{fontenc}
  \usepackage[utf8]{inputenc}
  \usepackage{textcomp} % provide euro and other symbols
\else % if luatex or xetex
  \usepackage{unicode-math} % this also loads fontspec
  \defaultfontfeatures{Scale=MatchLowercase}
  \defaultfontfeatures[\rmfamily]{Ligatures=TeX,Scale=1}
\fi
\usepackage{lmodern}
\ifPDFTeX\else
  % xetex/luatex font selection
  \setmainfont[]{Times New Roman}
  \setmonofont[]{Menlo}
  \setmathfont[]{STIX Two Math}
  \ifXeTeX
    \usepackage{xeCJK}
    \setCJKmainfont[]{Songti SC}
  \fi
  \ifLuaTeX
    \usepackage[]{luatexja-fontspec}
    \setmainjfont[]{Songti SC}
  \fi
\fi
% Use upquote if available, for straight quotes in verbatim environments
\IfFileExists{upquote.sty}{\usepackage{upquote}}{}
\IfFileExists{microtype.sty}{% use microtype if available
  \usepackage[]{microtype}
  \UseMicrotypeSet[protrusion]{basicmath} % disable protrusion for tt fonts
}{}
\setlength{\emergencystretch}{3em} % prevent overfull lines
\providecommand{\tightlist}{%
  \setlength{\itemsep}{0pt}\setlength{\parskip}{0pt}}
\usepackage[a4paper,top=24mm,bottom=25mm,left=23mm,right=23mm,headheight=15pt]{geometry}
\usepackage{amsmath,amssymb,mathtools,bm}
\usepackage{booktabs,longtable,array}
\usepackage{graphicx}
\usepackage{float}
\usepackage{caption}
\usepackage{enumitem}
\usepackage{fancyhdr}
\usepackage{titlesec}
\usepackage{mdframed}
\usepackage{xcolor}
\usepackage{xurl}
\usepackage{hyperref}
\graphicspath{{../../}}

\definecolor{PaperBlue}{HTML}{000000}
\definecolor{PaperRule}{HTML}{B7C6D5}
\definecolor{PaperShade}{HTML}{F4F7FA}
\definecolor{PaperText}{HTML}{000000}

\hypersetup{
  colorlinks=true,
  linkcolor=PaperBlue,
  urlcolor=PaperBlue,
  citecolor=PaperBlue,
  pdfborder={0 0 0}
}

\color{PaperText}
\linespread{1.08}
\setlength{\parindent}{1.5em}
\setlength{\parskip}{0.25em}
\setlength{\emergencystretch}{3em}
\setlist{itemsep=0.2em,topsep=0.45em}
\setcounter{secnumdepth}{3}
\setcounter{tocdepth}{3}

\titleformat{\section}
  {\Large\bfseries\color{PaperBlue}}
  {\thesection}{0.7em}{}
\titleformat{\subsection}
  {\large\bfseries\color{PaperText}}
  {\thesubsection}{0.65em}{}
\titleformat{\subsubsection}
  {\normalsize\bfseries\color{PaperText}}
  {\thesubsubsection}{0.6em}{}
\titlespacing*{\section}{0pt}{2.2ex plus 0.5ex}{0.9ex}
\titlespacing*{\subsection}{0pt}{1.8ex plus 0.4ex}{0.7ex}
\titlespacing*{\subsubsection}{0pt}{1.5ex plus 0.3ex}{0.5ex}

\captionsetup{
  font=small,
  labelfont=bf,
  textfont=it,
  justification=centering,
  skip=6pt
}

\pagestyle{fancy}
\fancyhf{}
\fancyhead[L]{\small\nouppercase{\leftmark}}
\fancyhead[R]{\small Selene's Blog}
\fancyfoot[C]{\small\thepage}
\renewcommand{\headrulewidth}{0.4pt}
\renewcommand{\footrulewidth}{0pt}
\fancypagestyle{plain}{
  \fancyhf{}
  \fancyhead[R]{\small Selene's Blog}
  \fancyfoot[C]{\small\thepage}
  \renewcommand{\headrulewidth}{0.4pt}
}

\renewenvironment{quote}
  {\begin{mdframed}[
      linewidth=0.7pt,
      linecolor=PaperRule,
      backgroundcolor=PaperShade,
      leftline=true,
      rightline=false,
      topline=false,
      bottomline=false,
      innerleftmargin=10pt,
      innerrightmargin=8pt,
      innertopmargin=7pt,
      innerbottommargin=7pt,
      skipabove=8pt,
      skipbelow=10pt
    ]\itshape}
  {\end{mdframed}}

\newenvironment{theorembox}[1]
  {\par\medskip\noindent\textbf{Theorem (#1).}\quad}
  {\par\medskip}

\newenvironment{lemmabox}[1]
  {\par\medskip\noindent\textbf{Lemma #1.}\quad}
  {\par\medskip}

\newenvironment{proofbox}
  {\par\medskip\noindent\textit{Proof.}\quad}
  {\hfill$\square$\par\medskip}

\AtBeginDocument{\allowdisplaybreaks[2]}
\usepackage{bookmark}
\IfFileExists{xurl.sty}{\usepackage{xurl}}{} % add URL line breaks if available
\urlstyle{same}
% fallback for those not using the hyperref driver hyperxmp:
\makeatletter
\@ifundefined{xmpquote}{\newcommand{\xmpquote}[1]{#1}}{}
\makeatother
\hypersetup{
  pdftitle={CS224N \textbar{} Assignment 2},
  pdfauthor={Selene},
  colorlinks=true,
  linkcolor={black},
  filecolor={Maroon},
  citecolor={Blue},
  urlcolor={black},
  pdfcreator={LaTeX via pandoc}}

\title{CS224N \textbar{} Assignment 2}
\author{Selene}
\date{August 25, 2026}

\begin{document}
\maketitle

\section{Problem 1}\label{problem-1}

Let \(y \in \mathbb{R}^{\lvert V\rvert}\) be the one-hot true
distribution, let \(\hat{y} \in \mathbb{R}^{\lvert V\rvert}\) be the
predicted distribution, and let

\[
\hat{y}_w
= P(O=w\mid C=c)
= \frac{\exp(u_w^\top v_c)}
{\sum_{j\in V}\exp(u_j^\top v_c)}.
\]

The matrix \(U\) contains the outside-word vectors as its columns:

\[
U = \begin{bmatrix}u_1 & u_2 & \cdots & u_{\lvert V\rvert}\end{bmatrix}.
\]

\subsection*{(a)}\label{a}
\addcontentsline{toc}{subsection}{(a)}

Since \(y\) is one-hot, \(y_o=1\) and \(y_w=0\) for every \(w\neq o\).
Therefore,

\[
-\sum_{w\in V}y_w\log \hat{y}_w
= -\log \hat{y}_o.
\]

\subsection*{(b)}\label{b}
\addcontentsline{toc}{subsection}{(b)}

\subsubsection*{\texorpdfstring{(i) Derivative with respect to
\(v_c\)}{(i) Derivative with respect to v\_c}}\label{i-derivative-with-respect-to-v_c}
\addcontentsline{toc}{subsubsection}{(i) Derivative with respect to
\(v_c\)}

First, rewrite the naive-softmax loss as

\[
J_{\text{naive-softmax}}(v_c,o,U)
= -\log \hat{y}_o
= -u_o^\top v_c+ \log\left(\sum_{w\in V}\exp(u_w^\top v_c)\right).
\]

Taking the derivative with respect to \(v_c\),

\[
\frac{\partial J}{\partial v_c}
= -u_o+ \frac{\sum_{w\in V}\exp(u_w^\top v_c)u_w}
{\sum_{j\in V}\exp(u_j^\top v_c)}
= -u_o + \sum_{w\in V}\hat{y}_w u_w.
\]

Because \(y\) is one-hot,

\[
u_o = Uy,\quad
\sum_{w\in V}\hat{y}_w u_w = U\hat{y}.
\]

Hence the vectorized result is

\[
\frac{\partial J}{\partial v_c}=U(\hat{y}-y)
\]

\subsubsection*{(ii) When the gradient is
zero}\label{ii-when-the-gradient-is-zero}
\addcontentsline{toc}{subsubsection}{(ii) When the gradient is zero}

The gradient is zero exactly when

\[
U(\hat{y}-y)=0
\]

Equivalently, \(\hat{y}-y\in\operatorname{Null}(U)\). In
particular,\(\hat{y}=y\) is sufficient for the gradient to be zero.

\subsubsection*{(iii) Interpretation of the two
terms}\label{iii-interpretation-of-the-two-terms}
\addcontentsline{toc}{subsubsection}{(iii) Interpretation of the two
terms}

Using a learning rate \(\eta>0\), gradient descent gives

\[
v_c
\leftarrow v_c-\eta\left(U\hat{y}-Uy\right)
=v_c+\eta u_o-\eta\sum_{w\in V}\hat{y}_w u_w.
\]

The \(+\eta u_o\) term moves the center-word vector \(v_c\) toward the
true outside-word vector \(u_o\), increasing their similarity. The
\(-\eta\sum_w\hat{y}_w u_w\) term moves \(v_c\) away from the
probability-weighted average of the predicted outside-word vectors, with
a larger correction for words to which the model currently assigns more
probability. Together, these changes increase the score of the true
outside word relative to the scores of the other words.

\subsection*{(c)}\label{c}
\addcontentsline{toc}{subsection}{(c)}

Suppose \(u_x=\alpha u_y\) for \(x\neq y\). If \(\alpha>0\), then

\[
\frac{u_x}{\lVert u_x\rVert_2}
=\frac{\alpha u_y}{\lVert \alpha u_y\rVert_2}
=\frac{u_y}{\lVert u_y\rVert_2}.
\]

Thus, L2 normalization makes two vectors that point in the same
direction identical, even if their original magnitudes differ. It
removes useful information when vector magnitude encodes information
relevant to the phrase classification---for example, when \(x\) and
\(y\) have similar semantic directions but different strengths or
polarities in the classifier. It does not remove useful information when
only vector direction matters, or when same-direction vectors with
different magnitudes should have the same effect on the final
classification.

\subsection*{(d)}\label{d}
\addcontentsline{toc}{subsection}{(d)}

Since \(U\) is formed by placing the outside-word vectors in columns,
its gradient has the same column structure:

\[
\frac{\partial J(v_c,o,U)}{\partial U}=
\begin{bmatrix}
\dfrac{\partial J(v_c,o,U)}{\partial u_1} &
\dfrac{\partial J(v_c,o,U)}{\partial u_2} &
\cdots &
\dfrac{\partial J(v_c,o,U)}{\partial u_{\lvert V\rvert}}
\end{bmatrix}
\]

\subsection*{(e)}\label{e}
\addcontentsline{toc}{subsection}{(e)}

The loss is

\[
J=-u_o^\top v_c+\log\left(\sum_{j\in V}\exp(u_j^\top v_c)\right).
\]

For the true outside word \(w=o\),

\[
\frac{\partial J}{\partial u_o}
=-v_c+\frac{\exp(u_o^\top v_c)}{\sum_{j\in V}\exp(u_j^\top v_c)}v_c
=(\hat{y}_o-1)v_c.
\]

For every \(w\neq o\),

\[
\frac{\partial J}{\partial u_w}
=\frac{\exp(u_w^\top v_c)}{\sum_{j\in V}\exp(u_j^\top v_c)}v_c
=\hat{y}_w v_c.
\]

Therefore,

\[
\frac{\partial J}{\partial u_w}
=(\hat{y}_w-y_w)v_c=
\begin{cases}
(\hat{y}_o-1)v_c, & w=o,\\
\hat{y}_w v_c, & w\neq o.
\end{cases}
\]

\subsection*{(f)}\label{f}
\addcontentsline{toc}{subsection}{(f)}

For \(f(x)=\max(\alpha x,x)\), where \(0<\alpha<1\),

\[
f'(x)=
\begin{cases}
1, & x>0,\\
\alpha, & x<0.
\end{cases}
\]

The derivative at \(x=0\) is not required.

\subsection*{(g)}\label{g}
\addcontentsline{toc}{subsection}{(g)}

Starting from

\[
\sigma(x)=\frac{1}{1+e^{-x}},
\]

we obtain

\[
\sigma'(x)
=\frac{e^{-x}}{(1+e^{-x})^2}
=\frac{1}{1+e^{-x}}\left(1-\frac{1}{1+e^{-x}}\right).
\]

Hence,

\[
\sigma'(x)=\sigma(x)\bigl(1-\sigma(x)\bigr)
\]

\section{Problem 2}\label{problem-2}

\subsection*{(a)}\label{a-1}
\addcontentsline{toc}{subsection}{(a)}

\subsubsection*{(i)}\label{i}
\addcontentsline{toc}{subsubsection}{(i)}

Momentum averages gradients over multiple minibatches, so random
variations in individual minibatch gradients tend to cancel out.
Directions that are consistently useful accumulate in m, making
parameter updates smoother and less noisy. This can reduce oscillation
and allow the model to converge more stably and quickly.

\subsubsection*{(ii)}\label{ii}
\addcontentsline{toc}{subsubsection}{(ii)}

Parameters whose gradients have historically been small receive
relatively larger updates, since their corresponding values in v are
smaller. Parameters with consistently large gradients receive smaller
updates. This adaptive scaling prevents large-gradient parameters from
changing too aggressively while still allowing small-gradient parameters
to make meaningful progress.

\end{document}
